\documentclass[a4paper,10pt]{article}
\usepackage[T2A]{fontenc}
\usepackage[utf8]{inputenc}

\usepackage[russian]{babel}

% Поиск в PDF
\usepackage{cmap}

\usepackage[a4paper,top=0.5cm,bottom=0.5cm,left=1cm,right=0.5cm,marginparwidth=0cm]{geometry}

% Красная строка для первого предложения в разделе
\usepackage{indentfirst}

% Равные промежутки между словами и предложениями
\frenchspacing

\usepackage{amsmath,amssymb,mathtools,amsthm}
\usepackage{icomma}

\usepackage{graphicx}

% Стиль нумерации формул
\renewcommand{\theequation}{\thesection.\arabic{equation}}

% Отображать номера только тех формул, на которые есть ссылка
\mathtoolsset{showonlyrefs=true}

% Более привычные греческие буквы
\renewcommand{\epsilon}{\varepsilon}
\renewcommand{\phi}{\varphi}
\renewcommand{\kappa}{\varkappa}

% Знаки сравнения с наклонёнными чёрточками
\renewcommand{\le}{\leqslant}
\renewcommand{\leq}{\leqslant}
\renewcommand{\ge}{\geqslant}
\renewcommand{\geq}{\geqslant}

\renewcommand{\emptyset}{\varnothing}

\newcommand{\ds}{\displaystyle}

% Обозначения числовых множеств
\newcommand{\N}{\mathbb{N}}
\newcommand{\Z}{\mathbb{Z}}
\newcommand{\Q}{\mathbb{Q}}
\newcommand{\R}{\mathbb{R}}

\DeclareMathOperator{\im}{Im}
\DeclareMathOperator{\re}{Re}

% Команда для переноса во встраиваемых в строку формулах. Взята из книги «Набор и вёрстка в системе
% LaTeX» С.М. Львовского.
\newcommand*{\hm}[1]{#1\nobreak\discretionary{}{\hbox{$\mathsurround=0pt #1$}}{}}

\usepackage[inline]{enumitem}

\sloppy

\usepackage{tikz}
\usetikzlibrary{decorations.pathreplacing}

\newcommand{\Section}[1]{{\medskip\Large \textbf{#1}\medskip}}
\newcommand{\Subsection}[1]{{\medskip\large \underline{\textbf{#1}}\medskip}}

\usepackage{float}

\newcommand{\alignv}{\vphantom{\left[\overrightarrow{OM}\right]}}
